
\subsubsection{Interpretation}

Our experiments validate the core hypothesis that joint DPO + attribute training is feasible through gradient-compatible multi-task optimization. The mean gradient angle of 78.5° provides quantitative proof that these tasks can be jointly optimized without catastrophic interference, extending multi-task learning theory to the LLM alignment domain where implicit reward modeling (DPO) and explicit user control (attributes) represent distinct but non-conflicting objectives.

The <120° threshold criterion offers a quantitative design principle for multi-objective LLM alignment: researchers can measure gradient angles to predict whether joint training will succeed before committing to expensive full-scale experiments. This principle generalizes to other alignment combinations where gradient angle analysis can guide architecture decisions.

The observed dual loss convergence (L\_DPO -5.8\%, L\_attr -21.3\%) demonstrates that joint optimization can achieve Pareto improvements where both objectives improve simultaneously. This challenges the common assumption that preference optimization and controllable generation require sequential stages to avoid catastrophic forgetting. While proof-of-concept experiments do not establish performance parity with standalone baselines, the absence of objective divergence provides strong evidence that full-scale joint training is viable.

The preference encoding finding (100\% probing accuracy) reveals that jointly trained models maintain task-specific representations despite multi-task pressure. This aligns with recent work on representation surgery for multi-task model merging \citep{Yang2024RepresentationSurgery}. Our gradient compatibility measurement (78.5° angle) provides the quantitative link explaining why this preservation occurs.

\subsubsection{Limitations}

\textbf{Proof-of-Concept Scale:} All experiments were conducted at approximately 1\% of planned training duration (100 vs 15,000 steps) due to computational constraints. This prevents us from claiming performance parity with standalone baselines. The observed performance gaps (0.5\% preference, 15\% steering) likely reflect incomplete convergence rather than fundamental incompatibility, as loss curves show continued decrease at training termination.

This limitation is acceptable because our research question addresses feasibility (can joint training work?) rather than optimization (does it match baselines?). The gradient compatibility finding (78.5° angle) is a step-level measurement that does not depend on full convergence, providing robust evidence of objective compatibility independent of training scale.

Future work: Full-scale 15,000-step training with loss weight ablation ($\alpha \in$ {0.5, 0.6, 0.7}) to optimize the preference-attribute tradeoff and close performance gaps.

\textbf{Synthetic Attribute Labels:} Attribute probing analysis used synthetic labels generated via random uniform distributions rather than real OpenAssistant annotations, yielding negative $R^{2}=-1.324$ that invalidates disentanglement measurement.

This limitation is acceptable because the preference encoding analysis functioned correctly (100\% accuracy), demonstrating that the probing methodology is sound. The negative $R^2$ is a clear failure signal allowing us to confidently discard attribute probing results while preserving preference findings.

Future work: Integrate real OpenAssistant attribute labels by mapping samples to HH-RLHF via shared prompts, enabling valid correlation measurement.

\textbf{Missing Sequential Baseline:} No sequential training baseline (DPO followed by attribute fine-tuning) was trained for comparison, preventing verification of emergent benefit claims.

This limitation is acceptable because feasibility is independently valuable for the alignment research community. Prior work has not demonstrated that DPO and attribute objectives can be jointly optimized---the default assumption is sequential training to avoid interference. Our gradient compatibility measurement provides quantitative evidence contradicting this assumption.

Future work: Train sequential baseline and compare to joint training on same held-out test set.
